% -------------------------------------------------------------------------------------
\section{推力分布的因子化}\label{sec:scet-thrust}

推力是 SCET$_{\rm I}$ 因子化的原型：观测量对红外安全（\cref{def:ir-safe}），末态只有出射的共线粒子，因而 Glauber 区不构成障碍；因子化定理在 SCET 中可以写到所有阶（[A]），而它的每一个单圈成分都可以独立计算并与全 QCD 比较（[R]，\cref{sec:scet-rg}）。

\paragraph{运动学。}
推力 $T=\max_{\hat{\bm t}}\sum_i|\bm p_i\cdot\hat{\bm t}|/\sum_i|\bm p_i|$\cite{Farhi1977}，$\tau\equiv1-T$；取质心系，$Q=\sqrt s$。推力轴把末态分为两个半球：$n$ 喷注所在的 R 半球（沿 $+z$）与 $\bar n$ 喷注所在的 L 半球。对无质量粒子 $\sum_i|\bm p_i|=Q$，推力轴平行于一个半球的总动量，故 $|\bm P_R|=|\bm P_L|=TQ/2$；又 $E_L+E_R=Q$ 且 $E_R^2-E_L^2=M_R^2-M_L^2$，即 $E_R-E_L=(M_R^2-M_L^2)/Q$。由 $M^2=E^2-|\bm P|^2$ 得 $M_L^2+M_R^2=E_L^2+E_R^2-\frac{T^2Q^2}{2}=\frac{Q^2}{2}(1-T^2)+\frac{(M_R^2-M_L^2)^2}{2Q^2}$，即\emph{精确}关系
\begin{equation}\label{eq:scet-thrust-hemi}
\frac{M_L^2+M_R^2}{Q^2}=\tau-\frac{\tau^2}{2}+\frac{(M_R^2-M_L^2)^2}{2Q^4},\qquad\text{故}\qquad \tau=\frac{M_L^2+M_R^2}{Q^2}+\order{\tau^2};
\end{equation}
% VF-SCET-15; checked on random 3- and 4-parton events (w18_s4, 5e-16)
三个部分子时 $M_L^2+M_R^2=Q^2\tau$ 精确成立。在领头幂次，$M_R^2=(p_n+p_{X_R})^2=p_n^2+\bar n\cdot p_n\,n\cdot p_{X_R}+\cdots=p_n^2+Q\,n\cdot p_{X_R}$（$\bar n\cdot p_n=Q[1+\order{\tau}]$），L 半球同理，于是
\begin{equation}\label{eq:scet-taudecomp}
\tau=\frac{p_R^2+p_L^2}{Q^2}+\frac kQ+\order{\tau^2},\qquad k\equiv n\cdot p_{X_R}+\bar n\cdot p_{X_L},
\end{equation}
其中 $p_{R}^2$、$p_L^2$ 是两个喷注的共线不变质量，$p_{X_R}$（$p_{X_L}$）是 R（L）半球中超软粒子的总动量。

\paragraph{硬、喷注与软函数。}
$\sigma_0=\frac{4\pi\alpha^2}{s}N_c\sum_qe_q^2\cdot\frac13$ 为 Born 截面。定义：(i) 硬函数 $H(Q^2,\mu)=|C_V(-Q^2-\ii0,\mu)|^2$，$C_V$ 是 QCD 矢量流 $\bar\psi\gamma^\mu\psi\to C_V\,[\bar\chi_n\gamma_\perp^\mu\chi_{\bar n}+(n\leftrightarrow\bar n)]$ 的 $\MSbar$ 匹配系数；(ii) 喷注函数由归一化明确的割定义
\begin{equation}\label{eq:scet-Jdef}
\sum_X(2\pi)^D\delta^D(p-p_X)\,\langle0|\chi^{(0)}_n|X\rangle\langle X|\bar\chi^{(0)}_n|0\rangle=2\pi\,\frac{\sla n}{2}\,\bar n\cdot p\;J(p^2,\mu)\;\mathbb 1_{\rm c}
\end{equation}
给出（$\mathbb 1_{\rm c}$ 为色单位矩阵，$\MSbar$ 重整化；树图层次 $P_n\sla pP_{\bar n}=\frac{\sla n}2\bar n\cdot p$ 给出 $J^{(0)}=\delta(p^2)$）；(iii) 推力软函数
\begin{equation}\label{eq:scet-Sdef}
S_T(k,\mu)=\frac{1}{N_c}\sum_{X}\Trc\Big[\langle0|\bar{\mathrm T}\big(Y_{\bar n}^\dagger Y_n\big)|X\rangle\langle X|\mathrm T\big(Y_n^\dagger Y_{\bar n}\big)|0\rangle\Big]\,
\delta\big(k-n\cdot p_{X_R}-\bar n\cdot p_{X_L}\big),
\end{equation}
树图层次 $S_T^{(0)}=\delta(k)$。

\begin{proposition}[硬匹配系数 \tagR]\label{prop:scet-hard-matching}
（$R_{\mathrm{PT}}$：在重整化微扰论中对红外正规化的部分子态严格。）（单圈。）记 $a\equiv\frac{C_F\as(\mu)}{4\pi}$。两喷注流的 $\MSbar$ 匹配系数与硬函数为
\begin{equation}\label{eq:scet-CV}
C_V(-q^2,\mu)=1+a\big[-L^2+3L-8+\zeta_2\big],\quad L\equiv\ln\frac{-q^2-\ii0}{\mu^2};\qquad
H(Q^2,\mu)=1+a\Big[-2L_Q^2+6L_Q-16+\frac{7\pi^2}{3}\Big],\quad L_Q\equiv\ln\frac{Q^2}{\mu^2},
\end{equation}
% VF-SCET-7, VF-SCET-8
并满足 $\frac{\dd C_V}{\dd\ln\mu}=[\Gamma_{\rm cusp}L+\gamma^V]C_V$、$\frac{\dd H}{\dd\ln\mu}=[2\Gamma_{\rm cusp}L_Q+2\gamma^H]H$，$\Gamma_0=4C_F$，$\gamma^H_0=\gamma^V_0=-6C_F$（展开约定 $\Gamma_{\rm cusp}=\sum_n(\frac{\as}{4\pi})^{n+1}\Gamma_n$，\cref{def:sudakov-cusp}）。这与\cite[附录~B]{BecherSchwartz2008}一致。（本命题只涉及部分子振幅。）
\uses{prop:sudakov-oneloop, lem:prelim-scaleless, def:sudakov-cusp}
\end{proposition}

\begin{proof}
无质量夸克的单圈裸形状因子精确于 $\varepsilon$（\cref{prop:sudakov-oneloop}；$\MSbar$ 中 $(4\pi\mu_0^2)^\varepsilon\to\mu^{2\varepsilon}\ee^{\varepsilon\gamma_E}$）为
\[
F^{(1)}=a\Big(\frac{\mu^2}{-q^2-\ii0}\Big)^{\varepsilon}\frac{\ee^{\varepsilon\gamma_E}\Gamma(1+\varepsilon)\Gamma^2(1-\varepsilon)}{\Gamma(1-2\varepsilon)}\Big[-\frac{2-\varepsilon+2\varepsilon^2}{\varepsilon^2(1-2\varepsilon)}\Big]
=a\Big[-\frac{2}{\varepsilon^2}+\frac{2L-3}{\varepsilon}-L^2+3L-8+\zeta_2+\order{\varepsilon}\Big],
\]
% VF-SCET-7 = VF-SUD-1
其中用了 $\ee^{\varepsilon\gamma_E}\Gamma(1+\varepsilon)\Gamma^2(1-\varepsilon)/\Gamma(1-2\varepsilon)=1-\frac{\zeta_2}{2}\varepsilon^2+\order{\varepsilon^3}$；所有极点都是红外的。另一方面，在壳 SCET 的单圈图（共线、超软）都无标度，按\cref{lem:prelim-scaleless}为零，故裸的 SCET 矩阵元等于树图，其紫外抵消项恰与 $F$ 的红外极点相同；匹配 $F=C_V\,\langle J_{\rm SCET}\rangle_{\rm ren}$ 于是给出 $C_V=F$ 的有限部分，即\cref{eq:scet-CV}的第一式。由 $\dd L/\dd\ln\mu=-2$，$\frac{\dd}{\dd\ln\mu}a(-L^2+3L)=a(4L-6)=\frac{\as}{4\pi}(\Gamma_0L+\gamma^V_0)$（$\dd a/\dd\ln\mu=\order{\as^2}$）。类时情形 $q^2=Q^2>0$，$L=L_Q-\ii\pi$，$\Real(L_Q-\ii\pi)^2=L_Q^2-\pi^2$，故 $|C_V|^2=1+2a[-L_Q^2+\pi^2+3L_Q-8+\frac{\pi^2}{6}]$，即 $H$ 的表达式，其 RG 方程由同样的求导给出。Becher--Schwartz 写 $H=1+\frac{\as}{4\pi}[-\frac12\Gamma_0L_Q^2-\gamma^H_0L_Q+c^H_1]$，$c^H_1=C_F(-16+\frac{7\pi^2}{3})$，与此相同。
\end{proof}

\begin{theorem}[推力分布的因子化 \tagA]\label{thm:scet-thrust}
对 $e^+e^-\to\gamma^*\to$ 强子，在 $\tau=1-T\ll1$ 的领头幂次，
\begin{equation}\label{eq:scet-thrust-fact}
\frac{1}{\sigma_0}\frac{\dd\sigma}{\dd\tau}=H(Q^2,\mu)\int\dd p_L^2\,\dd p_R^2\,\dd k\;J(p_L^2,\mu)\,J(p_R^2,\mu)\,S_T(k,\mu)\;
\delta\Big(\tau-\frac{p_L^2+p_R^2}{Q^2}-\frac kQ\Big)+\order{\tau^0},
\end{equation}
余项相对于 $1/\tau$ 型奇异项压低一个 $\tau$ 的幂次，在 $\tau\to0$ 可积；$H$、$J$、$S_T$ 由\cref{eq:scet-CV,eq:scet-Jdef,eq:scet-Sdef}定义。喷注函数必须用归一化明确的定义\cref{eq:scet-Jdef}：\cite[式~(2.6)]{BecherSchwartz2008}的写法 $J=\frac{1}{\bar n\cdot p}\frac1\pi\Imag[\ii\int\dd^4x\,\ee^{-\ii px}\langle0|\mathrm T\{\bar\chi_n(x)\frac{\bar{\sla n}}2\chi_n(0)\}|0\rangle]$ 若按字面计算（费米子编时的符号、对 Dirac 与色指标求迹），自由场给出 $-2N_c\,\delta(p^2)$，即其中隐含因子 $-\frac1{2N_c}\Tr$；与\cref{eq:scet-Jdef}等价、归一化明确的定义见\cite[式~(3)、(6)]{Becher:2006qw}。% E24 (convention hazard, Q3): kept in the statement
\par 非微扰地位：\cref{asm:intro-bridge}（MA1、MA2）。
\uses{thm:scet-decoupling, prop:scet-hard-matching, def:scet-zero-bin, prop:method-glauber-deformation, thm:ir-leading-regions, asm:intro-bridge}
\end{theorem}

\begin{proofsketch}
% errtype: E-S
(1)~\emph{硬匹配。}在 $\order{\lambda^0}$，$\bar\psi\gamma^\mu\psi\to C_V[\bar\chi_n\gamma_\perp^\mu\chi_{\bar n}+(n\leftrightarrow\bar n)]$：算符形式在树图层次由\cref{prop:scet-wilson}给出，共线规范不变性（\cref{prop:scet-gauge}）使它在各阶只能以 $\chi_n$、$\chi_{\bar n}$ 出现；$C_V$ 由在壳形状因子确定（\cref{prop:scet-hard-matching}）。
(2)~\emph{超软退耦。}由\cref{thm:scet-decoupling}与\cref{eq:scet-JY}，流化为 $C_V\bar\chi^{(0)}_nY_n^\dagger\gamma_\perp^\mu Y_{\bar n}\chi^{(0)}_{\bar n}$，而 $\mathcal L^{(0)}$ 分解为互不作用的 $n$、$\bar n$、超软三部分，态空间相应地分解为 $|X\rangle=|X_n\rangle|X_{\bar n}\rangle|X_{us}\rangle$。
(3)~\emph{测量的分解。}由\cref{eq:scet-taudecomp}，在领头幂次 $\tau$ 是三部分之和；共线粒子总在其喷注所在的半球。
(4)~\emph{求和。}对末态求和、对色指标缩并：喷注矩阵元正比于 $\mathbb 1_{\rm c}$，色单态流使软矩阵元以 $\frac1{N_c}\Trc$ 出现；轻子张量与 $\gamma_\perp^\mu$ 的 Dirac 迹给出 $\sigma_0$。于是截面化为\cref{eq:scet-thrust-fact}的卷积。
(5)~\emph{零箱。}由\cref{def:scet-zero-bin}(c)，推力的零箱积分无标度，减除后的喷注函数等于未减除的，与 $S_T$ 不重复计数。
(6)~\emph{Glauber。}所有共线粒子都是出射的，喷注线吸收软动量时的极点位于同一半平面（\cref{lem:coll-final-poles}），围道可以离开 Glauber 区（\cref{prop:method-glauber-deformation}(a)），故不需要 Glauber 算符。
(7)~\emph{区域。}硬、两个共线与超软模式正是 $e^+e^-\to$ 强子在测量 $\tau$ 时的领头区域（\cref{thm:ir-leading-regions}，以 $\lambda^2\sim\tau$ 计数）。
完整论证见\cite{Bauer:2003di,Fleming:2007qr,Schwartz:2007ib,Fleming:2007xt,BecherSchwartz2008,Bauer:2008dt}。
\end{proofsketch}

\paragraph{地位与缺口}
\cref{thm:scet-thrust}是在以硬、两个共线、超软为模式的 SCET 内对所有阶成立的论证（[A]）。最弱的一步是模式完备，即\cref{asm:intro-bridge}(Br1) 与积分形式的幂次计数\cref{asm:appN-envelope}：(7) 只给出 E-S 型的标度陈述，余项 $\order{\tau}$ 不是一个对所有阶一致的积分界。对 $e^+e^-$ 不存在 Glauber 钳夹，这一步依赖的是围道变形论证（\cref{prop:method-glauber-deformation,lem:coll-final-poles}），而不是 Rothstein--Stewart；后者只可作为一般机制引用：迭代的 Glauber 交换要么给出相位，要么可以吸收进软或共线 Wilson 线的方向，“因而在单举截面中相消”\cite{RothsteinStewart2016}。在 Collins 减除的语言中，推力的减除结构是\cref{prop:appN-thrust}（[R, R$_{\rm cond}$]），其前提\cref{asm:appN-collins}、\cref{asm:appN-cutsum}、\cref{asm:appN-ward-product}决定了图形路线的 [A]；图形路线中 Ward 恒等式这一步还牵涉第~9~章所命名的缺口（\cref{asm:appN-brst-nonlocal}，[C]），SCET 路线以有效理论的规范对称性代替它。非微扰地，当 $Q\tau\sim\Lambda$ 时软函数成为非微扰的形状函数\cite{Korchemsky:1994is,Korchemsky:1999kt,Bauer:2003di}；在尾部 $\Lambda/Q\ll\tau$，主要的强子化效应是分布的平移 $\propto\Lambda/Q$，作为在不同事件形状之间普适的陈述它属于 [C]（第~21~章）。文献：\cite{Bauer:2003di,Lee:2006nr,Fleming:2007qr,Schwartz:2007ib,Fleming:2007xt,BecherSchwartz2008,Bauer:2008dt}。升级途径：对 $e^+e^-$ 事件形状把\cref{asm:appN-envelope}与\cref{asm:appN-collins}在所有阶以积分范数形式证明。

\Cref{fig:dijet-factorization}画出了\cref{eq:scet-thrust-fact}的结构（图中 $s_n$、$s_{\bar n}$ 即 $p_R^2$、$p_L^2$，$S$ 即 $S_T$）。

%%% BEGIN fig:dijet-factorization
%%% (placeholder: replaced verbatim from figblocks.tex at assembly)
%%% END fig:dijet-factorization
